\section{Einleitung}

\gls{ai} und insbesondere \gls{llm} erreichen bei vielen Aufgaben hohe Leistung, ihr internes Verhalten bleibt aber oft undurchsichtig \cite{vaswani2017attention, radford2018improving}. Die \gls{mechanisticinterpretability} untersucht, wie und warum eine Ausgabe im Modell entsteht \cite{nanda2023comprehensive, elhage2021mathematical}.

Als Phänomen wird die \gls{sentimentanalysis} untersucht.
Eingesetzt wird dafür das offene Modell Pythia-410M von EleutherAI \cite{biderman2023pythia}.

Daraus ergeben sich drei Hypothesen zur Verarbeitung von Sentiment im Modell. Jede Hypothese ist mit einem konkreten Kriterium versehen, das angibt, wann sie als bestätigt gilt.

\begin{description}[style=nextline, leftmargin=0pt, font=\bfseries, itemsep=0.4em]
  \item[H1 (klare Stimmungswörter)] Eindeutig polare Wörter wie „wonderful“ oder „delicious“ werden früh und direkt am Token verarbeitet. Bestätigt, wenn die Polarität bereits in frühen bis mittleren Schichten ablesbar ist \emph{und} ein Eingriff am Stimmungswort die Vorhersage wiederherstellt.
  \item[H2 (Ironie und Sarkasmus)] Das Modell erkennt keine Ironie, sondern verarbeitet nur die wörtliche Oberflächen-Polarität. Bestätigt, wenn das Modell durchgehend der wörtlichen Lesart folgt und die gemeinte ironische Polarität an keiner Schicht linear ablesbar ist.
  \item[H3 (Lokalisierbarkeit)] Die kausal verantwortliche Verarbeitung lässt sich auf eine bestimmte Schicht und Token-Position eingrenzen, statt über das gesamte Netz verteilt zu sein. Bestätigt, wenn eine scharf begrenzte Stelle das korrekte Verhalten nahezu vollständig wiederherstellt.
\end{description}

